% ═══════════════════════════════════════════════════════════════
% AB-05a: Schwingungen
% Physik Klasse 10 MR
% Erkenntnisgeführt mit Versuchsprotokollen
% ═══════════════════════════════════════════════════════════════

\documentclass[11pt, a4paper]{article}

\usepackage[utf8]{inputenc}
\usepackage[T1]{fontenc}
\usepackage[ngerman]{babel}

\usepackage{helvet}
\renewcommand{\familydefault}{\sfdefault}

\usepackage{microtype}
\usepackage[margin=1.5cm, top=2cm, bottom=1.5cm]{geometry}
\usepackage{enumitem}
\usepackage{tabularx}
\usepackage{booktabs}
\usepackage{array}
\usepackage{tikz}
\usetikzlibrary{calc}

\usepackage{amsmath, amssymb}
\usepackage{siunitx}
\sisetup{locale = DE, per-mode = symbol}

\usepackage{fancyhdr}
\usepackage{lastpage}
\usepackage{needspace}

\usepackage{xcolor}
\usepackage{tcolorbox}
\tcbuselibrary{skins}

% ═══════════════════════════════════════════════════════════════
% FARBEN
% ═══════════════════════════════════════════════════════════════

\definecolor{physik}{HTML}{2c5aa0}
\definecolor{hellgrau}{HTML}{F5F5F5}
\definecolor{mittelgrau}{HTML}{888888}
\definecolor{rahmen}{HTML}{CCCCCC}
\definecolor{protokoll}{HTML}{fff3e0}
\definecolor{protokollrand}{HTML}{b35c00}

\colorlet{fachfarbe}{physik}

% ═══════════════════════════════════════════════════════════════
% VARIABLEN
% ═══════════════════════════════════════════════════════════════

\newcommand{\fach}{Physik}
\newcommand{\thema}{Schwingungen}
\newcommand{\klasse}{10 MR}
\newcommand{\stunde}{05a}

% ═══════════════════════════════════════════════════════════════
% KOPF- UND FUSSZEILE
% ═══════════════════════════════════════════════════════════════

\pagestyle{fancy}
\fancyhf{}
\renewcommand{\headrulewidth}{0.4pt}
\renewcommand{\footrulewidth}{0pt}
\lhead{\small \fach{} Klasse \klasse{} | \thema}
\rhead{\small Name: \underline{\hspace{3.5cm}}}
\cfoot{\small\textcolor{mittelgrau}{Seite \thepage{} von \pageref{LastPage}}}

% ═══════════════════════════════════════════════════════════════
% BOXEN
% ═══════════════════════════════════════════════════════════════

\newtcolorbox{merkkasten}[1][]{
    colback=hellgrau,
    colframe=hellgrau,
    coltitle=black,
    colbacktitle=hellgrau,
    boxrule=0pt,
    arc=0pt,
    left=8pt, right=8pt, top=8pt, bottom=6pt,
    borderline north={2pt}{0pt}{fachfarbe},
    before skip=6pt, after skip=6pt,
    fonttitle=\bfseries,
    #1
}

\newtcolorbox{protokollkasten}[1][]{
    colback=protokoll,
    colframe=protokoll,
    coltitle=black,
    colbacktitle=protokoll,
    boxrule=0pt,
    arc=0pt,
    left=8pt, right=8pt, top=8pt, bottom=6pt,
    borderline west={3pt}{0pt}{protokollrand},
    before skip=6pt, after skip=6pt,
    fonttitle=\bfseries,
    #1
}

\newtcolorbox{formelkasten}{
    colback=white,
    colframe=fachfarbe,
    boxrule=1.5pt,
    arc=0pt,
    left=8pt, right=8pt, top=6pt, bottom=6pt,
    before skip=4pt, after skip=4pt
}

% ═══════════════════════════════════════════════════════════════
% BEFEHLE
% ═══════════════════════════════════════════════════════════════

\newcommand{\luecke}[1]{\underline{\hspace{#1}}}
\newcommand{\punktebox}[1]{\hfill\small\textcolor{mittelgrau}{[\_\_/#1]}}

\newcommand{\aufgabe}[1]{%
    \needspace{4\baselineskip}%
    \vspace{10pt}%
    \noindent\textbf{\textcolor{fachfarbe}{Aufgabe #1}}%
    \vspace{4pt}%
    \nopagebreak[4]%
}

\newcommand{\operator}[1]{\textbf{\textcolor{fachfarbe}{#1}}}

\newlist{teil}{enumerate}{2}
\setlist[teil,1]{label=\alph*), leftmargin=1.8em, itemsep=3pt, parsep=0pt, topsep=3pt}

\newcommand{\antwortzeilen}[1]{%
    \par\vspace{4pt}%
    \foreach \n in {1,...,#1}{%
        \noindent\rule{\linewidth}{0.3pt}\vspace{6pt}%
    }%
}

\setlength{\parindent}{0pt}
\setlength{\parskip}{4pt}

% ═══════════════════════════════════════════════════════════════
% DOKUMENT
% ═══════════════════════════════════════════════════════════════

\begin{document}

% ═══════════════════════════════════════════════════════════════
% HEADER
% ═══════════════════════════════════════════════════════════════

\begin{center}
    {\Large\bfseries\textcolor{fachfarbe}{\thema}}\\[0.2em]
    {\small Arbeitsblatt | \fach{} Klasse \klasse{} | Stunde \stunde}
\end{center}
\vspace{-0.3em}
\hrule
\vspace{0.6em}

% ═══════════════════════════════════════════════════════════════
% K1: DEFINITION SCHWINGUNG
% ═══════════════════════════════════════════════════════════════

\begin{merkkasten}[title={\textbf{K1: Was ist eine Schwingung?}}]
Eine \textbf{Schwingung} ist eine \luecke{4cm} Bewegung um eine \luecke{3cm}.

\vspace{0.3em}
\textbf{Beispiele:} \luecke{2.5cm}, \luecke{2.5cm}, \luecke{2.5cm}, \luecke{2.5cm}
\end{merkkasten}

% ═══════════════════════════════════════════════════════════════
% P1: VERSUCHSPROTOKOLL FADENPENDEL
% ═══════════════════════════════════════════════════════════════

\begin{protokollkasten}[title={\textbf{P1: Versuchsprotokoll -- Fadenpendel}}]

\textbf{Forschungsfrage:} Wovon hängt die Periodendauer $T$ beim Fadenpendel ab?

\vspace{0.5em}
\textbf{1. Hypothese} (Vermutung VOR dem Experiment):

Ich vermute, dass $T$ von \luecke{3cm} abhängt, weil \luecke{5cm}.

\vspace{0.5em}
\textbf{2. Durchführung} -- Messtabelle:

\vspace{0.3em}
\begin{tabular}{|c|c|c|c|c|}
\hline
\textbf{Messung} & \textbf{Länge $L$ (cm)} & \textbf{Masse $m$ (g)} & \textbf{Amplitude $A$ (cm)} & \textbf{Periodendauer $T$ (s)} \\
\hline
1 & & & & \\
\hline
2 & & & & \\
\hline
3 & & & & \\
\hline
4 & & & & \\
\hline
\end{tabular}

\vspace{0.5em}
\textbf{3. Beobachtung:}

Wenn ich die Länge $L$ vergrößere, dann \luecke{4cm} die Periodendauer $T$.

Wenn ich die Masse $m$ vergrößere, dann \luecke{4cm} die Periodendauer $T$.

Wenn ich die Amplitude $A$ vergrößere, dann \luecke{4cm} die Periodendauer $T$.

\vspace{0.5em}
\textbf{4. Erkenntnis:}

$T$ hängt ab von: \luecke{3cm}

$T$ hängt \textbf{nicht} ab von: \luecke{3cm} und \luecke{3cm}
\end{protokollkasten}

% ═══════════════════════════════════════════════════════════════
% P2: VERSUCHSPROTOKOLL FEDERPENDEL
% ═══════════════════════════════════════════════════════════════

\begin{protokollkasten}[title={\textbf{P2: Versuchsprotokoll -- Federpendel}}]

\textbf{Forschungsfrage:} Wovon hängt die Periodendauer $T$ beim Federpendel ab?

\vspace{0.5em}
\textbf{1. Hypothese} (Vermutung VOR dem Experiment):

Ich vermute, dass $T$ von \luecke{3cm} abhängt, weil \luecke{5cm}.

\vspace{0.5em}
\textbf{2. Durchführung} -- Messtabelle:

\vspace{0.3em}
\begin{tabular}{|c|c|c|c|c|}
\hline
\textbf{Messung} & \textbf{Federhärte $D$ (N/m)} & \textbf{Masse $m$ (g)} & \textbf{Amplitude $A$ (cm)} & \textbf{Periodendauer $T$ (s)} \\
\hline
1 & & & & \\
\hline
2 & & & & \\
\hline
3 & & & & \\
\hline
4 & & & & \\
\hline
\end{tabular}

\vspace{0.5em}
\textbf{3. Beobachtung:}

Wenn ich die Federhärte $D$ vergrößere, dann \luecke{4cm} die Periodendauer $T$.

Wenn ich die Masse $m$ vergrößere, dann \luecke{4cm} die Periodendauer $T$.

Wenn ich die Amplitude $A$ vergrößere, dann \luecke{4cm} die Periodendauer $T$.

\vspace{0.5em}
\textbf{4. Erkenntnis:}

$T$ hängt ab von: \luecke{3cm} und \luecke{3cm}

$T$ hängt \textbf{nicht} ab von: \luecke{3cm}
\end{protokollkasten}

\newpage

% ═══════════════════════════════════════════════════════════════
% K2: KENNGROSSEN
% ═══════════════════════════════════════════════════════════════

\begin{merkkasten}[title={\textbf{K2: Die drei Kenngrößen einer Schwingung}}]
\begin{tabularx}{\linewidth}{|l|c|X|c|}
\hline
\textbf{Größe} & \textbf{Symbol} & \textbf{Bedeutung} & \textbf{Einheit} \\
\hline
Amplitude & $A$ & \luecke{5cm} & m (oder cm) \\
\hline
Periodendauer & $T$ & \luecke{5cm} & \luecke{2cm} \\
\hline
Frequenz & $f$ & \luecke{5cm} & \luecke{2cm} \\
\hline
\end{tabularx}
\end{merkkasten}

% ═══════════════════════════════════════════════════════════════
% K3: FORMEL
% ═══════════════════════════════════════════════════════════════

\begin{formelkasten}
\begin{center}
\textbf{K3: Die Formel}

\vspace{0.5em}
{\Large $f = \dfrac{1}{T}$} \qquad bzw. \qquad {\Large $T = \dfrac{1}{f}$}

\vspace{0.5em}
\textbf{Bedeutung:} Je \luecke{2cm} die Frequenz, desto \luecke{2cm} die Periodendauer.
\end{center}
\end{formelkasten}

% ═══════════════════════════════════════════════════════════════
% D1: y(t)-DIAGRAMM
% ═══════════════════════════════════════════════════════════════

\begin{merkkasten}[title={\textbf{D1: Das $y(t)$-Diagramm}}]
\textbf{Zeichne} ein $y(t)$-Diagramm mit $A = 10\,\text{cm}$ und $T = 2\,\text{s}$:

\vspace{0.3em}
\begin{center}
\begin{tikzpicture}[scale=0.85]
    % Achsen
    \draw[->] (0,0) -- (9,0) node[right] {$t$/s};
    \draw[->] (0,-2.5) -- (0,2.5) node[above] {$y$/cm};

    % Y-Skala
    \foreach \y in {-10,-5,5,10} {
        \draw (-0.1,\y/5) -- (0.1,\y/5);
        \node[left] at (-0.15,\y/5) {\small \y};
    }
    \node[left] at (-0.15,0) {\small 0};

    % X-Skala
    \foreach \x in {1,2,3,4,5,6,7,8} {
        \draw (\x,-0.1) -- (\x,0.1);
        \node[below] at (\x,-0.15) {\small \x};
    }

    % Gitter (leicht)
    \foreach \x in {1,2,3,4,5,6,7,8} {
        \draw[gray!30] (\x,-2.2) -- (\x,2.2);
    }
    \foreach \y in {-2,-1,1,2} {
        \draw[gray!30] (0,\y) -- (8.5,\y);
    }
\end{tikzpicture}
\end{center}

\vspace{0.3em}
\textbf{So liest du das Diagramm:}
\begin{itemize}[leftmargin=2em, itemsep=0pt]
    \item $y$-Achse: \luecke{3cm} (in cm)
    \item $t$-Achse: \luecke{3cm} (in s)
    \item Amplitude $A$: Abstand von $y=0$ zum \luecke{2cm}
    \item Periodendauer $T$: Abstand zwischen zwei \luecke{2cm}
\end{itemize}
\end{merkkasten}

% ═══════════════════════════════════════════════════════════════
% AUFGABEN
% ═══════════════════════════════════════════════════════════════

\aufgabe{Ü1: Frequenz und Periodendauer berechnen} \punktebox{6}

\operator{Berechne} die fehlende Größe.

\begin{teil}
    \item $T = 4\,$s. Berechne $f$.

    \vspace{0.3em}
    $f = $ \luecke{6cm} $= $ \luecke{2cm}
    \vspace{0.5em}

    \item $f = 2\,$Hz. Berechne $T$.

    \vspace{0.3em}
    $T = $ \luecke{6cm} $= $ \luecke{2cm}
    \vspace{0.5em}

    \item $T = 0{,}5\,$s. Berechne $f$.

    \vspace{0.3em}
    $f = $ \luecke{6cm} $= $ \luecke{2cm}
\end{teil}

% ---

\aufgabe{Ü2: Periodendauer aus Messung berechnen} \punktebox{4}

Ein Pendel schwingt 10-mal in 20 Sekunden.

\begin{teil}
    \item \operator{Berechne} die Periodendauer $T$.

    \vspace{0.3em}
    $T = $ \luecke{6cm} $= $ \luecke{2cm}

    \item \operator{Berechne} die Frequenz $f$.

    \vspace{0.3em}
    $f = $ \luecke{6cm} $= $ \luecke{2cm}
\end{teil}

% ---

\aufgabe{Ü3: Diagramm ablesen} \punktebox{6}

\operator{Lies} die Werte aus dem Diagramm ab und \operator{berechne} die Frequenz.

\begin{center}
\begin{tikzpicture}[scale=0.75]
    % Achsen
    \draw[->] (0,0) -- (8,0) node[right] {$t$/s};
    \draw[->] (0,-2.2) -- (0,2.2) node[above] {$y$/cm};

    % Y-Skala
    \node[left] at (-0.1,2) {\small 20};
    \node[left] at (-0.1,0) {\small 0};
    \node[left] at (-0.1,-2) {\small -20};
    \draw (-0.1,2) -- (0.1,2);
    \draw (-0.1,-2) -- (0.1,-2);

    % X-Skala
    \foreach \x in {1,2,3,4,5,6,7} {
        \draw (\x,-0.1) -- (\x,0.1);
        \node[below] at (\x,-0.15) {\small \x};
    }

    % Sinuskurve: A=20cm, T=4s
    \draw[physik, thick, domain=0:7.5, samples=100]
        plot (\x, {2*sin(\x*90)});
\end{tikzpicture}
\end{center}

Amplitude $A$ = \luecke{2cm} \hfill Periodendauer $T$ = \luecke{2cm} \hfill Frequenz $f$ = \luecke{2cm}

% ---

\aufgabe{Ü4: Transfer -- Herzfrequenz} \punktebox{4}

Ein Herz schlägt in Ruhe 72-mal pro Minute.

\begin{teil}
    \item \operator{Berechne} die Frequenz $f$ in Hz.

    \vspace{0.3em}
    Hinweis: 1 Minute = 60 Sekunden

    \vspace{0.3em}
    $f = $ \luecke{6cm} $= $ \luecke{2cm}

    \item \operator{Berechne} die Periodendauer $T$ (Zeit für einen Herzschlag).

    \vspace{0.3em}
    $T = $ \luecke{6cm} $= $ \luecke{2cm}
\end{teil}

% ═══════════════════════════════════════════════════════════════
% PUNKTE
% ═══════════════════════════════════════════════════════════════

\vfill
\hrule
\vspace{0.4em}

\hfill\textbf{Punkte:} \luecke{1.2cm} / 20

\end{document}
